\documentclass[a4paper,fontset=none]{ctexart}

\usepackage{indentfirst}
\setlength{\parindent}{0em}
\usepackage{autobreak}
\allowdisplaybreaks
\usepackage{parskip}
\usepackage{geometry}
\geometry{top=3cm,bottom=3cm,left=2cm,right=2cm}
\usepackage{anyfontsize}

\usepackage{fontspec}
\setmainfont{STIX Two Text}
\usepackage{unicode-math}
\unimathsetup{mathrm=sym,mathbf=sym,mathsf=sym,bold-style=ISO}
\setmathfont{STIX Two Math}
\usepackage{physics2}
\usephysicsmodule{ab,op.legacy}
\usepackage{fixdif,derivative}
\newcommand{\um}{\upmu\mathrm{m}}
\newcommand{\ang}{\text{\r{A}}}
\AtBeginDocument{
    \let\uglyepsilon\epsilon
    \renewcommand\epsilon{\varepsilon}
}

\setCJKmainfont{Source Han Serif CN}[BoldFont=Source Han Sans CN Medium,ItalicFont=FZKai-Z03]

\title{\textbf{星际介质天文学作业}}
\author{张植竣}
\date{2024年5月10日}

\begin{document}

\maketitle

\begin{enumerate}
    \item[31.1 (a)]
    根据化学反应定律，对于基元反应$a\mathrm{A} + b\mathrm{B} \to c\mathrm{C} + d\mathrm{D}$，其反应速率为：
    \[
        R = \frac{1}{a} \odv{n(\mathrm{A})}{t} = \frac{1}{b} \odv{n(\mathrm{B})}{t} = -\frac{1}{c} \odv{n(\mathrm{C})}{t} = -\frac{1}{d} \odv{n(\mathrm{D})}{t} = k n(\mathrm{A})^a n(\mathrm{B}^b)
    \]
    故对于生成$\mathrm{H}^-$的反应$\mathrm{H} + e^- \to \mathrm{H}^- +h\nu$，反应速率为：
    \[
        R_1 = \odv{n(\mathrm{H}^-)}{t} = k_{\mathrm{ra}} n(\mathrm{H}) n(e^-)
    \]
    对于消耗$\mathrm{H}^-$的反应$\mathrm{H}^- + \mathrm{H} \to \mathrm{H}_2 + e^-$，反应速率为：
    \[
        R_2 = -\odv{n(\mathrm{H}^-)}{t} = k_{\mathrm{ad}} n(\mathrm{H}^-) n(\mathrm{H})
    \]
    对于消耗$\mathrm{H}^-$的反应$\mathrm{H}^- + h\nu \to \mathrm{H} + e^-$，反应速率为：
    \[
    R_3 = -\odv{n(\mathrm{H}^-)}{t} = \zeta_{\mathrm{pd}} n(\mathrm{H}^-)
    \]
    体系平衡稳定时，生成$\mathrm{H}^-$的速率与消耗$\mathrm{H}^-$的速率相等，此时：
    \[
        R_1 = R_2 + R_3
    \]
    即：
    \[
        k_{\mathrm{ra}} n(\mathrm{H}) n(e^-) = k_{\mathrm{ad}} n(\mathrm{H}^-) n(\mathrm{H}) + \zeta_{\mathrm{pd}} n(\mathrm{H}^-)
    \]
    则有：
    \begin{align*}
        \frac{n(\mathrm{H}^-)}{n(\mathrm{H})}
        &= \frac{k_{\mathrm{ra}} n(e^-)}{k_{\mathrm{ad}} n(\mathrm{H}) + \zeta_{\mathrm{pd}}} \\
        &= \frac{1.9 \times 10^{-16} T_2^{0.67} \times 0.02}{1.3 \times 10^{-9} \times 30 + 2.4 \times 10^{-7}} \\
        &= 1.362 \times 10^{-11} T_2^{0.67}
    \end{align*}

    \medskip
    \item[31.1 (b)]
    参与$\mathrm{H}^-$的反应$\mathrm{H}^- + \mathrm{H} \to \mathrm{H}_2 + e^-$反应的$\mathrm{H}^-$比例为：
    \begin{align*}
        \frac{R_2}{R_2 + R_3}
        &= \frac{k_{\mathrm{ad}} n(\mathrm{H}^-) n(\mathrm{H})}{k_{\mathrm{ad}} n(\mathrm{H}^-) n(\mathrm{H}) + \zeta_{\mathrm{pd}} n(\mathrm{H}^-)} \\
        &= \frac{k_{\mathrm{ad}} n(\mathrm{H})}{k_{\mathrm{ad}} n(\mathrm{H}) + \zeta_{\mathrm{pd}}} \\
        &= \frac{1.3 \times 10^{-9} \times 30}{1.3 \times 10^{-9} \times 30 + 2.4 \times 10^{-7}} \\
        &= 0.14
    \end{align*}

    \medskip
    \item[31.1 (c)]
    由题意：
    \begin{align*}
        R_{\mathrm{H}^-}
        &= \frac{k_{\mathrm{ad}} n(\mathrm{H}^-) n(\mathrm{H})}{n(\mathrm{H}) n(\mathrm{H})} \\
        &= k_{\mathrm{ad}} \frac{n(\mathrm{H}^-)}{n(\mathrm{H})} \\
        &= 1.3 \times 10^{-9} \times 1.36 \times 10^{-11} T_2^{0.67} \\
        &= 1.77 \times 10^{-20} T_2^{0.67}\,\mathrm{cm^3/s}
    \end{align*}
    这个速率比尘埃催化生成$\mathrm{H}_2$的速率（由(31.11)式给出）：
    \[
        R_{\mathrm{gr}} = 7.3 \times 10^{17} T_2^{1/2} \ab<\epsilon_{\mathrm{gr}}> \Sigma_{-21}\,\mathrm{cm^3/s} \approx 3.65 \times 10^{-17} T_2^{1/2}\,\mathrm{cm^3/s}
    \]
    小了3个数量级。

\end{enumerate}

\end{document}